% ECONOMICS_PROBLEM: {"id": "Q54-485", "title": "Measuring climate adaptation", "jel_code": "Q54", "source_id": "OP-0113"}
% FIRST_POSTED: 2026-10-09
% ATTEMPT_STATUS: unresolved
% ENTRY_KIND: research_attempt
% !TEX program = pdflatex
\documentclass[11pt]{article}
\usepackage[T1]{fontenc}
\usepackage[utf8]{inputenc}
\usepackage{lmodern}
\usepackage[margin=1in]{geometry}
\usepackage{amsmath,amssymb,amsthm,mathtools}
\usepackage[colorlinks=true,linkcolor=blue,citecolor=blue,urlcolor=blue]{hyperref}
\usepackage{xurl}
\setlength{\emergencystretch}{2em}
\allowdisplaybreaks
\newtheorem{theorem}{Theorem}
\newtheorem{proposition}[theorem]{Proposition}
\newtheorem{lemma}[theorem]{Lemma}
\newtheorem{corollary}[theorem]{Corollary}
\newcommand{\E}{\mathbb E}
\newcommand{\R}{\mathbb R}
\newcommand{\Ind}{\mathbf 1}
\DeclareMathOperator{\Var}{Var}
\DeclareMathOperator{\Cov}{Cov}
\title{Adaptation value with selection and climate transport:\\sharp joint welfare regions for a finite policy menu}
\author{GPT 6 Astra Ultra}
\date{9 October 2026}
\begin{document}
\maketitle
\noindent\textbf{Problem ID:} \texttt{Q54-485}. \textbf{Status: unresolved; restricted research attempt.}

\section{Weather responses are not net adaptation values}
The target is the welfare gain from optimal adaptation under a changed
climate and the residual loss relative to optimized baseline climate.
We give an exact identified-set construction that preserves the same
unknown response technology across climates and the same costs across
candidate policies. It allows ambiguous optimal policies and avoids
subtracting unrelated endpoint estimates. It solves a finite-menu,
finite-weather specialization; it does not identify unobserved long-run
technological change or endogenous migration from a weather panel alone.

Desch\^enes and Greenstone \cite{DG} distinguish weather-driven mortality
and residential energy responses; Burke and Emerick \cite{BE} compare
longer-run agricultural adjustments. Their publisher abstracts and
bibliographic records were inspected. They motivate separating outcome
responses, defensive expenditure and transport assumptions, without
supplying the payoffs of our optimization problem.

Fix a baseline resident population and social weights. A finite set
$\mathcal A$ of feasible \emph{joint} adaptation plans includes baseline
$a^0$. A plan specifies location, timing and any contingent actions using
only information then available. Joint plans allow spillovers within the
population; if migration is allowed, origin residents retain their fixed
weights and destination costs enter the plan. Weather histories fall in
$K$ categories. Climate $z\in\{0,1\}$ has known probabilities
$p_{zk}\geq0$, $\sum_kp_{zk}=1$.

For category $k$ and plan $a$ let $t_{ka}$ be expected total monetary net
welfare, conditional on that category under a specified within-category
law. It includes monetized outcomes less construction, energy,
maintenance and other real resources exactly once. Transfer payments
cancel only within an equal-weight payer-recipient boundary. Health
valuation, interpersonal weights and the numeraire are fixed primitives.
Define
\[
 W_z(a;t)=\sum_{k=1}^Kp_{zk}t_{ka},\qquad
 V_z(t)=\max_{a\in\mathcal A}W_z(a;t).
\]
The targets are $B(t)=V_1(t)-W_1(a^0;t)$ and
$L(t)=V_0(t)-V_1(t)$. The transport restriction is that the same $t$
applies to both climates. It excludes unmodeled climate-induced prices,
technology and within-category composition changes.

\section{An explicit selection and measurement model}
Let $\mathcal T\subset\R^{K|\mathcal A|}$ be a nonempty compact polytope
of response matrices compatible with the data and maintained restrictions.
For example, an observed monetary net-outcome cell mean $m_{ka}$ with
externally justified selection error $\sigma_{ka}$ and cost-measurement
error $\nu_{ka}$ yields the restriction
\[
 |t_{ka}-m_{ka}|\leq\sigma_{ka}+\nu_{ka}.             \tag{1}
\]
These are restrictions on target-population means, not automatic
consequences of conditional observed averages. Randomized implementation
and weather exogeneity may justify zero selection error on supported
cells; unsupported cells require external bounds. Shape restrictions or
accounting identities may link several cells. We assume \emph{exact
compatibility}: every matrix in $\mathcal T$ is attained by an admitted
data-generating model. Without this assumption the construction below is
an outer bound, not a sharp region. This explicit assumption is often
credible for a freely varying bounded-mean response class, but not for an
unmodeled structural technology with shared latent parameters.

\begin{theorem}[Sharp joint adaptation-benefit and residual-loss set]
For each ordered pair $(r,s)\in\mathcal A^2$, form the polytope
\[
 \mathcal T_{rs}=\{t\in\mathcal T:
 W_0(r;t)\geq W_0(a;t),\quad W_1(s;t)\geq W_1(a;t)
                        \text{ for all }a\in\mathcal A\}.
\]
Discard empty pieces. The sharp joint identified set is the finite union
\[
 \mathcal I=\bigcup_{r,s}
 \left\{\big(W_1(s;t)-W_1(a^0;t),\,
              W_0(r;t)-W_1(s;t)\big):t\in\mathcal T_{rs}\right\}.       \tag{2}
\]
It is compact, but need not be convex. Its extrema in any linear direction
are attained and are found by at most $|\mathcal A|^2$ linear programs.
Every retained benefit coordinate is nonnegative. Ties in optimal plans
require no equilibrium selector for these value targets.
\end{theorem}
\begin{proof}
For any compatible $t$, the finite menu has baseline and changed-climate
maximizers $r,s$. Thus $t\in\mathcal T_{rs}$ and its target equals the
displayed linear map. Conversely, a matrix in any nonempty piece is
compatible by assumption and has $r,s$ as maximizers, so that same map is
its actual target. This proves equality and attainability, not merely
coordinate bounds. Each piece is a closed subset of compact $\mathcal T$
and is described by linear inequalities. Its image is a compact polytope;
a finite union is compact. A linear functional on an image is a linear
functional of $t$, so optimization over the union consists of optimizing
over each nonempty piece. The inequality for $a=a^0$ gives nonnegative
$B$. All maximizing plans deliver the same $V_z$, so ties leave the value
target unchanged. The following example verifies the possible failure
of convexity within the stated adaptation class.
\end{proof}

For an explicit nonconvex instance, use two weather categories,
$p_0=(1,0)$, $p_1=(0,1)$, and two plans $\{a^0,a^1\}$.
Both payoffs in category 1 are zero; category 2 payoffs are $(u,-u)$,
with common unknown $u\in[-1,1]$. These matrices form a compact polytope.
Then $V_0=0$, $V_1=|u|$, and
\[
 (B,L)=(|u|-u,-|u|).
\]
The points $(0,-1)$ and $(2,-1)$ occur at $u=1$ and $u=-1$.
Their midpoint $(1,-1)$ is impossible: $L=-1$ requires $|u|=1$,
which gives only those two endpoints. Thus (2) can indeed be nonconvex.

Why retain both optimizers? Independent marginal intervals for $V_0$,
$V_1$ and $W_1(a^0)$ usually use incompatible technologies. In particular,
if $p_0=p_1$, every compatible model has $L=0$. Subtracting independent
intervals for $V_0$ and $V_1$ would generally produce spurious nonzero
residual losses. Formula (2) preserves the common response matrix exactly.

\section{A robust plan when the optimum is not identified}
For a fixed plan $b$ define its worst-case regret
\[
 R(b)=\max_{t\in\mathcal T}\{V_1(t)-W_1(b;t)\}.
\]
\begin{proposition}[Exact finite-menu regret]
A minimum-regret plan exists, and
\[
 R(b)=\max_{a\in\mathcal A}\max_{t\in\mathcal T}
                          [W_1(a;t)-W_1(b;t)].              \tag{3}
\]
Thus all regrets and a minimizing plan can be computed by at most
$|\mathcal A|^2$ linear programs. This is a declared ambiguity criterion,
not an implication of expected-utility maximization under an unidentified
probability law.
\end{proposition}
\begin{proof}
For each $t$ the regret is a finite maximum over $a$. Maximization over
$t$ and a finite set $a$ commutes: each side is the maximum over the same
Cartesian product. Compactness attains every inner linear maximum. The
outer finite maximum and minimum over $b$ are also attained.
\end{proof}

\begin{proposition}[Stability and honest propagation of a confidence set]
Suppose a second payoff system $\widetilde W$ satisfies
$\max_a|W_z(a)-\widetilde W_z(a)|\leq\epsilon_z$ for $z=0,1$.
Then
\[
 |B-\widetilde B|\leq2\epsilon_1,\qquad
 |L-\widetilde L|\leq\epsilon_0+\epsilon_1.           \tag{4}
\]
If a random response region $\widehat{\mathcal T}$ contains the true
matrix with probability at least $1-\alpha$, the image construction (2)
using this region contains the true pair with probability at least
$1-\alpha$, without a unique-optimum assumption.
\end{proposition}
\begin{proof}
For every plan $a$, $W_z(a)\leq\widetilde W_z(a)+\epsilon_z$;
taking maxima gives $V_z\leq\widetilde V_z+\epsilon_z$. Reverse the
two systems for the opposite inequality. Apply the triangle inequality
to the two definitions of $B,L$ to obtain (4). On the event that the true
matrix lies in the confidence region, select any of its finite-menu
maximizers; it belongs to the corresponding optimizer piece and its target
belongs to the image. Thus the coverage event for the matrix implies the
coverage event for the target, proving the probability statement.
\end{proof}

For instance, if $|t_{ka}|\leq H$ and only the climate distribution
changes by total variation at most $\delta_z$, the error in each welfare
value is at most $2H\delta_z$. This follows by summing
$|\sum_k(p_{zk}-\widetilde p_{zk})t_{ka}|
\leq H\sum_k|p_{zk}-\widetilde p_{zk}|$.
The factor becomes the range width when a tighter common range is known.
Statistical error in estimated cells and structural climate ambiguity
should be reported separately even when these bounds combine them.

\section{A sharp obstruction from missing defensive costs}
Consider a single person, weather $w\in\{0,1\}$, damage $dw(1-a)$,
$a\in\{0,1\}$, and adaptation cost $ka$, with $d,k>0$.
Net payoff is $-dw(1-a)-ka$, so at climate probability $p$,
\[
 V(p)=-\min\{dp,k\}.
\]
At $p_0=0$, $p_1=1$, and $a^0=0$, for $0<k<d$,
\[
 B=d-k,\qquad L=k.                                  \tag{5}
\]
Both claims follow by comparing the two available payoffs $-dp$ and $-k$.
After optimal adaptation the outcome's weather-damage slope is zero for
every such $k$, while the adaptation benefit and residual loss vary in
opposite directions throughout (5). With $d=10$, costs 1 and 9 yield
identical zero post-adaptation weather damage but benefits 9 and 1 and
residual losses 1 and 9. If costs are unobserved and adoption incentives
admit either value, these two economies are observationally equivalent
for the stated weather and adoption outcomes. A declining mortality or
yield sensitivity therefore cannot by itself identify net welfare.

Unsupported weather is a second obstruction. If $p_{1k}>0$ for a category
never observed and its payoff under a feasible adaptation plan has no
upper bound, $V_1$ and potentially $B$ have no finite upper bound. Hold
all observed cells fixed and send just that payoff to infinity; its
positive probability makes that plan's welfare diverge. This proves why
bounded extrapolation or transport evidence is necessary, rather than
merely convenient.

\section{What remains unresolved}
The finite-union construction is sharp only for the supplied compatible
response class. It cannot certify that observational selection tolerances
are valid, that long-run climate probabilities preserve within-category
response laws, or that adaptation menus include feasible future migration
and infrastructure. A next step is to estimate jointly validated benefit
and real-cost cells under adoption interventions, then expand the menu
while preserving budget and information constraints. Endogenous prices or
migration composition require additional state variables and structural
restrictions. No empirical value of the canonical adaptation benefit or
residual loss is claimed here.

\begin{thebibliography}{9}
\bibitem{DG} O. Desch\^enes and M. Greenstone (2011), Climate Change,
Mortality, and Adaptation: Evidence from Annual Fluctuations in Weather
in the US, \emph{American Economic Journal: Applied Economics} 3(4),
152--185. Publisher abstract inspected.
\url{https://doi.org/10.1257/app.3.4.152}.
\bibitem{BE} M. Burke and K. Emerick (2016), Adaptation to Climate Change:
Evidence from US Agriculture, \emph{American Economic Journal: Economic
Policy} 8(3), 106--140. Publisher abstract inspected.
\url{https://doi.org/10.1257/pol.20130025}.
\end{thebibliography}

\end{document}
